\SourceLeaf{32}
\fancyhead[L]{\small OBSERVATIONES COMETAE SECUNDI A. MDCCCXIII}
\renewcommand{\Article}[1]{\par\medskip\Needspace{4\baselineskip}\pdfbookmark[1]{Art. #1}{cometae-art-#1}\centerline{\textbf{#1.}}\nobreak\smallskip}
\pdfbookmark[0]{Observationes Cometae secundi A. MDCCCXIII}{cometae}
\begin{center}
{\large OBSERVATIONES}\par\medskip
{\Large COMETAE SECUNDI A. MDCCCXIII}\par\medskip
IN OBSERVATORIO GOTTINGENSI FACTAE\par\medskip
ADIECTIS\par\medskip
NONNULLIS ADNOTATIONIBUS CIRCA CALCULUM ORBITARUM PARABOLICARUM.
\end{center}

Cometam a collega amicissimo cl. \textsc{Harding} d. 3 Aprilis huius anni in constellatione Tauri Poniatovii detectum ipse in specula nostra observare coepi inde a d. 7 Aprilis. Ecce determinationes, quas ope micrometri circularis telescopio decempedali adaptati obtinere licuit:
\begin{center}\small\setlength{\tabcolsep}{4pt}
\begin{tabular}{r|r|r|r l}
1813&T. med. Gott.&Asc. R. app.&\multicolumn{2}{c}{Declin. app.}\\\hline
Apr. 7&$13^{\mathrm h}\;12^{\mathrm m}\;2^{\mathrm s}$&$271^0\;7'\;19''3$&$5^0\;34'\;36''7$&Bor.\\
9&$13\;35\;40$&$270\;10\;33.5$&$4\;11\;3.4$&—\\
11&$13\;17\;43$&$269\;1\;19.9$&$2\;33\;0.7$&—\\
14&$13\;7\;36$&$266\;44\;5.5$&$0\;33\;0.8$&Austr.\\
21&$14\;23\;0$&$256\;39\;19.3$&$12\;57\;56.0$&—
\end{tabular}\end{center}

Postea ad quadrantem muralem a clar. \textsc{Harding} observationes factae sunt hae:
\begin{center}\small\setlength{\tabcolsep}{4pt}
\begin{tabular}{r|r|r|r l}
Apr. 21&$15^{\mathrm h}\;7^{\mathrm m}\;21^{\mathrm s}$&$256^0\;34'\;19''6$&$13^0\;2'\;26''5$&Austr.\\
24&$14\;22\;50$&$248\;23\;21$&$21\;45\;2$&\\
25&$14\;4\;21$&$244\;44\;42$&$25\;10\;42$&
\end{tabular}\end{center}

D. 24 et 25 Aprilis cometa nudis quoque oculis valde conspicuus fuit. Noctibus insequentibus coelum nubibus tectum, moxque rapidus cometae ad austrum descensus finem observationibus imposuit.

\SourceLeaf{33}
Superfluum videtur, hocce loco elementa parabolica repetere, quae statim ab initio e tribus primis observationibus ipse deduxeram. Etenim horum accuratius expoliendorum curam demandavi calculatori valde exercitato atque perito, Doctori \textsc{Gerling}, cui sequentia elementa correcta debemus, cunctis observationibus tum nostris tum iis quas transmiserat clar. \textsc{Olbers}, quantum licuit adaptata.
\begin{center}\fontsize{9}{11}\selectfont\setlength{\tabcolsep}{1.5pt}
\begin{tabular}{@{}l l r@{}}
Logarithmus distantiae in perihelio&. . . . . . . . . .&$0.0849212$\\
\multicolumn{3}{@{}p{\linewidth}@{}}{Tempus transitus per perihelium, ad meridianum Gottingensem}\\
\multicolumn{3}{r@{}}{1813 Mai. $19.44507$}\\
Longitudo perihelii&. . . . . . . . . . . . . .&$197^0\;43'\;7''7$\\
Longitudo nodi ascendentis&. . . . , . . . . . .&$42\;40\;15.2$\\
Inclinatio orbitae&. . . . . . . . . . . . . . .&$81\;2\;11.8$\\
\multicolumn{3}{@{}l}{Motus \textit{retrogradus}.}
\end{tabular}\end{center}

\Needspace{11\baselineskip}
Observationes clar. \textsc{Olbers} fuerunt hae:
\begin{center}\small\setlength{\tabcolsep}{4pt}
\begin{tabular}{r|r|r|r l}
1813&T. med. Brem.&Asc. R. app.&\multicolumn{2}{c}{Declin. app.}\\\hline
Apr. 14&$13^{\mathrm h}\;31^{\mathrm m}\;4^{\mathrm s}$&$266^0\;42'\;51''2$&$0^0\;34'\;22''8$&Austr.\\
15&$12\;14\;29$&$265\;48\;47.9$&$1\;46\;4.5$&\\
19&$11\;38\;0$&$260\;40\;39.1$&$8\;15\;23.7$&\\
21&$12\;0\;35$&$256\;51\;59.3$&$12\;42\;54.3$&\\
24&$11\;58\;38$&$248\;43\;57.7$&$21\;25\;9.8$&\\
25&$11\;41\;30$&$245\;8\;18.0$&$24\;49\;2.4$&\\
&$12\;5\;38$&$245\;4\;3.0$&$24\;54\;16.4$&
\end{tabular}\end{center}

Accessit observatio clar. \textsc{Bouvard} in observatorio imp. Parisino facta:
\begin{center}\small\setlength{\tabcolsep}{4pt}
\begin{tabular}{r|r|r|r l}
1813&T. med. Paris.&Asc. R. app.&\multicolumn{2}{c}{Declin. app.}\\\hline
Apr. 13&$16^{\mathrm h}\;22^{\mathrm m}\;2^{\mathrm s}$&$267^0\;27'\;18''$&$0^0\;24'\;46''$&Bor.
\end{tabular}\end{center}

Consensus elementorum supra traditorum cum singulis hisce observationibus e calculis clar. \textsc{Gerling} ita se habet:
\SourceLeaf{34}
\begin{center}\small\setlength{\tabcolsep}{6pt}
\begin{tabular}{r|r|r|l}
&\multicolumn{2}{c|}{Differentia}&\\
&Asc. R.&Declin.&Observator\\\hline
Apr. 7&$+3''8$&$+8''5$&Gauss\\
9&$+2.0$&$+34.3$&Gauss\\
11&$-5.3$&$-17.7$&Gauss\\
13&$-1.6$&$-0.4$&Bouvard\\
14&$-7.4$&$-28.6$&Gauss\\
&$+2.7$&$-8.4$&Olbers\\
15&$-0.9$&$+28.7$&Olbers\\
19&$-25.1$&$+103.9$&Olbers\\
21&$-56.6$&$-59.2$&Olbers\\
&$-30.1$&$-5.2$&Gauss\\
&$-22.8$&$-24.1$&Harding\\
24&$-45.2$&$-41.4$&Olbers\\
&$+0.4$&$-11.6$&Harding\\
25&$-23.3$&$-67.8$&Olbers\\
&$-9.4$&$-27.4$&Olbers\\
&$+9.7$&$+1.4$&Harding
\end{tabular}\end{center}

In hacce comparatione tum aberrationis. tum parallaxis ratio rite est habita.
\begin{center}\rule{.25\linewidth}{.3pt}\end{center}

Liceat his addere quaedam calculi compendia, quibus saepius in determinatione prima orbitae parabolicae secundum methodum clar. \textsc{Olbers} commode usus sum, et per quae methodus ista iam per se tam expedita adhuc magis contrahi vel ad calculum numericum magis idonea reddi videtur. Referuntur ea ad computum radiorum vectorum, atque imprimis chordae inter locum primum et ultimum. Clar. \textsc{Olbers} ad hunc finem adhibet expressiones huius formae $\BareRoot{(f+g\rho+h\rho\rho)}$, atque coëfficientes $f$, $g$, $h$ per operationes satis quidem simplices determinat, ita tamen comparatas, ut in plerisque casibus praecisio sufficiens obtineri nequeat, nisi tabulis logarithmorum maioribus ad septem vel saltem ad sex figuras decimales constructis perficiantur. Illarum expressionum loco alias substitui, quae tum per se aliquanto magis commodae ad calculum videntur, tum eo quoque nomine se commendant, quod omnibus operationibus quinque decimales adeoque tabulae logarithmorum minores sufficiunt. Rei summa hisce momentis innititur. Sint
\SourceLeaf{35}
\begin{description}
\item[$\odot$, $\odot'$, $\odot''$] longitudines Solis in observatione prima, secunda, tertia
\item[$R$, $R'$, $R''$] distantiae solis a terra
\item[$\alpha$, $\alpha'$, $\alpha''$] longitudines atque
\item[$\ObservedBeta$, $\ObservedBeta'$, $\ObservedBeta''$] latitudines geocentricae cometae
\item[$r$, $r'$, $r''$] distantiae eiusdem a sole
\item[$\rho$, $\rho'$, $\rho''$] eiusdem distantiae curtatae a terra
\item[$t$, $t'$, $t''$] tempora observationum
\item[$k$] chorda inter cometae locum primum atque tertium
\item[$M=\frac{\rho''}{\rho}$]
\end{description}

His positis facile perspicitur haberi
\[
\begin{aligned}{}
[1]\quad r&=\BareRoot{[(\rho\cos\alpha-R\cos\odot)^2}\\
&\qquad +(\rho\sin\alpha-R\sin\odot)^2+\rho\rho\operatorname{tang}\ObservedBeta^2]\\
[2]\quad r''&=\BareRoot{[(M\rho\cos\alpha''-R''\cos\odot'')^2}\\
&\qquad +(M\rho\sin\alpha''-R''\sin\odot'')^2\\
&\qquad +MM\rho\rho\operatorname{tang}{\ObservedBeta''}^2]\\
[3]\quad k&=\BareRoot{[(M\rho\cos\alpha''-\rho\cos\alpha-R''\cos\odot''+R\cos\odot)^2}\\
&\qquad +(M\rho\sin\alpha''-\rho\sin\alpha-R''\sin\odot''+R\sin\odot)^2\\
&\qquad +(M\rho\operatorname{tang}\ObservedBeta''-\rho\operatorname{tang}\ObservedBeta)^2]
\end{aligned}
\]
Aequationes 1, 2 induunt formam hancce
\[
\begin{aligned}
r&=\BareRoot{\left(\frac{\rho\rho}{\cos\ObservedBeta^2}-2\rho R\cos(\alpha-\odot'')+RR\right)}\\
r''&=\BareRoot{\left(\frac{MM\rho\rho}{\cos{\ObservedBeta''}^2}-2M\rho R''\cos(\alpha''-\odot'')+R''R''\right)}
\end{aligned}
\]
Statuendo itaque
\[
\begin{aligned}
\cos\ObservedBeta\cos(\alpha-\odot)&=\cos\psi, & R\sin\psi&=B\\
\cos\ObservedBeta''\cos(\alpha''-\odot'')&=\cos\psi'', & R''\sin\psi''&=B''
\end{aligned}
\]
habebimus
\[
\begin{aligned}
r&=\BareRoot{\left[\left(\frac{\rho}{\cos\ObservedBeta}-R\cos\psi\right)^2+BB\right]}\\
r''&=\BareRoot{\left[\left(\frac{M\rho}{\cos\ObservedBeta''}-R''\cos\psi''\right)^2+B''B''\right]}
\end{aligned}
\]

Determinando porro quinque quantitates auxiliares $g$, $G$, $h$, $H$, $\zeta$, ita ut habeatur
\SourceLeaf{36}
\[
\begin{aligned}
R''\cos\odot''-R\cos\odot&=g\cos G\\
R''\sin\odot''-R\sin\odot&=g\sin G\\
M\cos\alpha''-\cos\alpha&=h\cos\zeta\cos H\\
M\sin\alpha''-\sin\alpha&=h\cos\zeta\sin H\\
M\operatorname{tang}\ObservedBeta''-\operatorname{tang}\ObservedBeta&=h\sin\zeta
\end{aligned}
\]
formula 3 transit in sequentem
\[
\begin{aligned}
k&=\BareRoot{[(\rho h\cos\zeta\cos H-g\cos G)^2}\\
&\qquad +(\rho h\cos\zeta\sin H-g\sin G)^2+\rho\rho hh\sin\zeta^2]\\
&=\BareRoot{(\rho\rho hh-2\rho hg\cos\zeta\cos(G-H)+gg)}
\end{aligned}
\]
Quodsi itaque statuimus
\[
\cos\zeta\cos(G-H)=\cos\varphi,\quad g\sin\varphi=A
\]
erit
\[
k=\BareRoot{[(\rho h-g\cos\varphi)^2+AA]}
\]
aut si insuper statuimus $\rho h-g\cos\varphi=u$,
\[
k=\BareRoot{(uu+AA)}
\]

Gratum fore censemus pluribus lectoribus, si non modo complexum omnium operationum ad has transformationes pertinentium rite ordinatum hic denuo sistamus, sed insuper omnes reliquas operationes simul adiiciamus, ita ut omnia, quae ad calculum primum orbitae parabolicae requiruntur, hic iuncta inveniantur. Simul haec praecepta per numeros a cometa nostro desumtos illustrabimus. Eligimus itaque observationes nostras dierum 7, 14, 21 Aprilis, quarum reductio sequentia \textit{data} administrat:
\[
\begin{array}{rcrrr@{\qquad}rcrrr}
t&=&\multicolumn{3}{l}{7.55002}&&&&&\\
t'&=&\multicolumn{3}{l}{14.54694}&&&&&\\
t''&=&\multicolumn{3}{l}{21.59931}&&&&&\\
\alpha&=&271^0&16'&38''&\ObservedBeta&=&+29^0&2'&0''\\
\alpha'&=&266&27&22&\ObservedBeta'&=&+22&52&18\\
\alpha''&=&256&48&8&\ObservedBeta''&=&+9&53&12\\
\odot&=&17&47&41&\log R&=&\multicolumn{3}{l}{0.00091}\\
\odot'&=&24&38&45&\log R'&=&\multicolumn{3}{l}{0.00175}\\
\odot''&=&31&31&25&\log R''&=&\multicolumn{3}{l}{0.00260}
\end{array}
\]

\SourceLeaf{37}
I. Operatio \textit{prima} consistit in determinatione valoris approximati ipsius $M$, quem suppeditat formula
\[
M=\frac{t''-t'}{t'-t}.\frac{\operatorname{tang}\ObservedBeta'\sin(\alpha-\odot')-\operatorname{tang}\ObservedBeta\sin(\alpha'-\odot')}{\operatorname{tang}\ObservedBeta''\sin(\alpha'-\odot')-\operatorname{tang}\ObservedBeta'\sin(\alpha''-\odot')}
\]
In exemplo nostro invenitur $\log M=9.75799$.

II. Tunc determinare oportet quantitates $g$, $G$, $h$, $H$, $\zeta$ per formulas sequentes, quas supra traditis aequivalentes sed ad calculum commodiores esse patet.
\[
\begin{aligned}
R''\cos(\odot''-\odot)-R&=g\cos(G-\odot)\\
R''\sin(\odot''-\odot)&=g\sin(G-\odot)\\
M-\cos(\alpha''-\alpha)&=h\cos\zeta\cos(H-\alpha'')\\
\sin(\alpha''-\alpha)&=h\cos\zeta\sin(H-\alpha'')\\
M\operatorname{tang}\ObservedBeta''-\operatorname{tang}\ObservedBeta&=h\sin\zeta
\end{aligned}
\]

In exemplo nostro provenit
\[
\begin{array}{rcrrr}
G&=&113^0&43'&57''\\
\log g&=&\multicolumn{3}{l}{9.38029}\\
H&=&109^0&5'&49''\\
\zeta&=&44&13&9\\
\log h&=&\multicolumn{3}{l}{9.81477}
\end{array}
\]

III. Postea statuemus
\[
\begin{aligned}
\cos\zeta\cos(G-H)&=\cos\varphi\\
\cos\ObservedBeta\cos(\alpha-\odot)&=\cos\psi\\
\cos\ObservedBeta''\cos(\alpha''-\odot'')&=\cos\psi''\\[5pt]
g\sin\varphi&=A\\
R\sin\psi&=B\\
R''\sin\psi''&=B''
\end{aligned}
\]
In hoc calculo si forte cosinus angulorum $\varphi$, $\psi$, $\psi''$ parum ab unitate diversi evadant, figuras sex vel adeo septem adhibere conveniet. Ceterum angulos $\varphi$, $\psi$, $\psi''$ ipsos gradibus minutis et secundis describere non necessarium est, quum sufficiat, statim a logarithmis cosinuum ad logarithmos sinuum transire.

\SourceLeaf{38}
In exemplo nostro fit
\[
\begin{aligned}
\log A&=9.22527\\
\log B&=9.98706\\
\log B''&=9.86038
\end{aligned}
\]

IV. Denique statuatur
\[
\begin{aligned}
h\cos\ObservedBeta&=b\\
\frac{h\cos\ObservedBeta''}{M}&=b''\\
g\cos\varphi-bR\cos\psi&=c\\
g\cos\varphi-b''R''\cos\psi''&=c''
\end{aligned}
\]
In exemplo nostro
\[
\begin{aligned}
\log b&=9.75645\\
\log b''&=0.05028\\
c&=+0.31365\\
c''&=+0.95443
\end{aligned}
\]

V. His ita praeparatis radii vectores $r$, $r''$ atque chorda $k$ pendent ab incognita $u$ sequenti modo
\[
\begin{aligned}
r&=\BareRoot{\left[\left(\frac{u+c}{b}\right)^2+BB\right]}\\
r''&=\BareRoot{\left[\left(\frac{u+c''}{b''}\right)^2+B''B''\right]}\\
k&=\BareRoot{(uu+AA)}
\end{aligned}
\]
quam incognitam ita per tentamina determinare oportet, ut satisfiat aequationi
\[
(r+r''+k)^{\frac32}-(r+r''-k)^{\frac32}=\frac{t''-t}{m}
\]
ubi notandum, $m$ designare tempus $9.6887401$ dierum atque esse
\[
\log m=0.9862673
\]
Quantitati $(r+r''-k)^{\frac32}$ signum $+$ praefigere oporteret, si motus heliocentricus cometae intra tempus $t''-t$ angulum $180^0$ superaret, sed hic casus in suppositionibus, quibus haecce prima orbitae determinatio innititur, occurrere nequit. Ceterum vix opus erit monere, calculum numericum pro $r$ ita instituendum esse, ut introducatur angulus auxiliaris $\theta$, talis ut sit
\SourceLeaf{39}
\[
\frac{bB}{u+c}=\operatorname{tang}\theta
\]
unde fieri $r=\frac{B}{\sin\theta}$, ac perinde pro $r''$ et $k$. Nec non quisque sponte videbit, in omnibus his operationibus percommode adhiberi posse tabulam nostram pro logarithmis summarum atque differentiarum immediate inveniendis.

In exemplo nostro habetur $\log\frac{t''-t}{m}=0.16139$, paucisque tentaminibus factis eruitur
\[
u=0.24388
\]

VI. Inventa quantitate $u$, habemus
\[
\rho=\frac{u+g\cos\varphi}{h},\quad\rho''=M\rho
\]
\begin{center}(in exemplo nostro $\log\rho=9.80364$, $\log\rho''=9.56163$).\end{center}

Operationes reliquae satis quidem sunt notae: sed ut omnia hic adsint, formulas reliquas quoque, quibus uti solemus, apponere visum est. Sint itaque
\begin{description}
\item[$\lambda$, $\lambda''$] longitudines heliocentricae cometae in observatione prima atque tertia
\item[$\ApproxBeta$, $\ApproxBeta''$] latitudines heliocentricae
\item[$\upsilon$, $\upsilon''$] longitudines in orbita
\item[$\AscendingNode$] longitudo nodi ascendentis
\item[$i$] inclinatio orbitae intra $0$ et $90^0$ accipienda, si secundum modum vulgarem motum directum retrogradumque distinguimus
\item[$\Pi$] longitudo perihelii
\item[$T$] tempus transitus per perihelium
\item[$q$] distantia in perihelio.
\end{description}

VII. Positiones heliocentricae invenientur per formulas
\[
\begin{aligned}
\rho\cos(\alpha-\odot)-R&=r\cos\ApproxBeta\cos(\lambda-\odot)\\
\rho\sin(\alpha-\odot)&=r\cos\ApproxBeta\sin(\lambda-\odot)\\
\rho\operatorname{tang}\ObservedBeta&=r\sin\ApproxBeta\\
\rho''\cos(\alpha''-\odot'')-R''&=r''\cos\ApproxBeta''\cos(\lambda''-\odot'')\\
\rho''\sin(\alpha''-\odot'')&=r''\cos\ApproxBeta''\sin(\lambda''-\odot'')\\
\rho''\operatorname{tang}\ObservedBeta''&=r''\sin\ApproxBeta''
\end{aligned}
\]

\SourceLeaf{40}
Consensus valorum pro radiis vectoribus $r$, $r''$ hinc prodeuntium cum iis, qui antea ex $u$ deducti fuerant, calculo confirmando inserviet. Motus directus erit vel retrogradus, prout $\lambda''$ maior evadit quam $\lambda$, vel minor.

In exemplo nostro invenimus
\[
\begin{aligned}
\lambda&=225^0\;4'\;22'',&\ApproxBeta&=+14^0\;51'\;39'',&\log r&=0.13896\\
\lambda''&=223\;6\;55,&\ApproxBeta''&=+2\;49\;28,&\log r''&=0.11068
\end{aligned}
\]

Cometae itaque motus est retrogradus.

VIII. Ad inveniendam longitudinem nodi atque inclinationem, adhibeantur formulae:
\[
\begin{aligned}
\pm\operatorname{tang}\ApproxBeta&=\operatorname{tang}i.\sin(\lambda-\AscendingNode)\\
\pm\frac{\operatorname{tang}\ApproxBeta''-\operatorname{tang}\ApproxBeta.\cos(\lambda''-\lambda)}{\sin(\lambda''-\lambda)}&=\operatorname{tang}i.\cos(\lambda-\AscendingNode)
\end{aligned}
\]
ubi signa superiora referuntur ad motum directum, inferiora ad retrogradum. Dein longitudines in orbita eruuntur per formulas
\[
\begin{aligned}
\frac{\operatorname{tang}(\lambda-\AscendingNode)}{\cos i\cdot}&=\operatorname{tang}(\upsilon-\AscendingNode)\\
\frac{\operatorname{tang}(\lambda''-\AscendingNode)}{\cos i}&=\operatorname{tang}(\upsilon''-\AscendingNode)
\end{aligned}
\]
ubi $\upsilon-\AscendingNode$, $\upsilon''-\AscendingNode$ in iisdem resp. quadrantibus accipere oportet, in quibus sunt $\lambda-\AscendingNode$, $\lambda''-\AscendingNode$.

Pro cometa nostro invenimus
\[
\begin{array}{rcrrr}
\AscendingNode&=&42^0&40'&8''\\
i&=&81&1&3\\
\upsilon&=&237&43&7\\
\upsilon''&=&225&31&32
\end{array}
\]

IX. Longitudinem perihelii atque distantiam in perihelio dant formulae
\[
\begin{gathered}
\frac{1}{\BareRoot r}=\frac{1}{\BareRoot q}.\cos\tfrac12(\upsilon-\Pi)\\
\frac{\operatorname{cotg}\tfrac12(\upsilon''-\upsilon)}{\BareRoot r}-\frac{1}{\sin\tfrac12(\upsilon''-\upsilon).\BareRoot{r''}}=\frac{1}{\BareRoot q}.\sin\tfrac12(\upsilon-\Pi)
\end{gathered}
\]
Pro cometa nostro $\Pi=197^0\;37'\;51''$, $\log q=0.08469$.

\SourceLeaf{41}
X. Denique e tabula \textsc{Barkeri} desumantur motus medii anomaliis veris $\upsilon-\Pi$, $\upsilon''-\Pi$ vel $\Pi-\upsilon$, $\Pi-\upsilon''$ respondentes, qui sint $M$, $M''$. Tunc erit
\[
T=t\mp Mnq^{\frac32}=t''\mp M''nq^{\frac32}
\]
ubi signa superiora valent, si in motu directo $\upsilon>\Pi$, $\upsilon''>\Pi$ vel in motu retrogrado $\upsilon<\Pi$, $\upsilon''<\Pi$; inferiora in casibus oppositis. Quantitas $n$ est constans, ipsiusque logarithmus $=0.0398723$. Consensus duorum valorum ipsius $T$ secundam calculi confirmationem subministrat.

In exemplo nostro invenimus
\[
\begin{aligned}
T&=49.518\\
T&=49.517
\end{aligned}
\]
ita ut pro tempore transitus per perihelium adoptari possit Maii d. $19.5175$.

Quodsi ex hisce elementis locus geocentricus pro tempore observationis mediae computatur, prodit longitudo $266^0\;27'\;15''$, latitudo $22^0\;52'\;18''$ bor., illa $7''$ ab observata discrepans, haec ex asse consentiens.
\par\medskip\begin{center}\rule{.25\linewidth}{.3pt}\end{center}
